%! AP Statistics — Unit 7, Lesson 7.3 — Homework
\documentclass[10pt]{article}
\usepackage{apstats-article}
\usepackage{apstats-boxes}

\begin{document}

\pageheader{Unit 7, Lesson 7.3}{Homework}
\namedateperiod

\vspace{0.1in}
\textbf{Directions:} For each problem, show your reasoning clearly.

\begin{enumerate}

  \item \textbf{Interpretation Practice.} A 90\% $t$-interval for the mean
        number of hours per week adults spend exercising is $(3.8,\ 5.6)$ hours,
        based on a random sample of 45 adults.

        \begin{enumerate}[label=\alph*.]
          \item Write a correct interpretation of the interval. \writelines{2}

          \item A fitness organization claims adults exercise a mean of 6 hours per
                week. Does the interval provide convincing evidence against this
                claim? Justify. \writelines{2}

          \item A health agency claims adults exercise fewer than 5 hours per week
                on average. Does the interval provide convincing evidence that the
                mean is less than 5 hours? Explain. \writelines{2}
        \end{enumerate}

  \item \textbf{Width Relationships.} A researcher constructs a 95\% $t$-interval
        for the mean price of textbooks at a university: $\bar x = \$87.40$,
        $s = \$24.00$, $n = 36$.

        \begin{enumerate}[label=\alph*.]
          \item Calculate the margin of error and the interval. Show work.
                \vspace{0.1in}

                $df = $ \blank{2cm}, $t^* = $ \blank{2.5cm}, $SE = $ \blank{2.5cm}

                $ME = $ \blank{3cm}, Interval: \blank{5cm} $= ($ \blank{2cm}$,$ \blank{2cm}$)$

          \item How would the margin of error change if $n = 144$? Compute it.
                \writelines{2}

          \item How would the interval change if the confidence level were raised to
                99\% (keeping $n = 36$)? Would you expect it to be wider or narrower?
                \writelines{2}
        \end{enumerate}

  \item \textbf{AP-Style Free Response.} A random sample of 22 employees at a
        large company reports a mean commute distance of 14.6 miles with a standard
        deviation of 8.3 miles.

        \begin{enumerate}[label=\alph*.]
          \item Construct and interpret a 95\% confidence interval for the true mean
                commute distance for all employees. Verify all conditions. \writelines{7}

                Interval: $= ($ \blank{2cm}$,$ \blank{2cm}$)$ miles

          \item The company claims the mean commute is 12 miles. Based on your interval,
                does the data provide convincing evidence that the mean commute exceeds
                12 miles? \writelines{2}
        \end{enumerate}

\end{enumerate}

\begin{reflectionbox}
\textbf{Reflection:} A classmate says ``A 99\% CI is always better than a 95\% CI
because it is more confident.'' Do you agree or disagree? What trade-off does the
student ignore?

\writelines{3}
\end{reflectionbox}

\end{document}
